The angular momentum of a particle is given by
\begin{align}
\bm{L}&=\bm{r} \times \bm{p} \nonumber \\
&=m \bm{r}\times \frac{d \bm{r}}{d t} \nonumber \\
&=m r \hat{\bm{r}} \times \left[ \dot{r} \hat{\bm{r}} + r \omega \hat{\bm{\theta}} \right] \nonumber \\
&=m \left[ r \dot{r}\hat{\bm{r}} \times \hat{\bm{r}} + r^2 \omega \hat{\bm{r}} \times\hat{\bm{\theta}} \right] \nonumber \\
&=m r^2 \omega \uveck \label{eq:AngularMomentum}
\end{align}

We can also write this as $\mathbf{L} = l \uveck$ were
\begin{equation}
l=mr^2\omega \label{eq:AngularMomentumScalar}
\end{equation}

\subsection{Conservation Of Angular Momentum}
We can show now that angular momentum is conserved under a central force by taking the derivative with respect to time.

\begin{equation}
\frac{d\mathbf{L}}{dt} = \frac{d l \uveck}{dt} = \frac{d l}{dt}\uveck
\end{equation}
\begin{align}
\frac{dl}{dt}&=m\frac{d}{dt}\left[r^2\omega\right]\nonumber \\
&=m\left[\frac{dr^2}{dt}\omega + r^2\frac{d\omega}{dt}\right]\nonumber \\
&=m\left[\frac{dr^2}{dr}\frac{dr}{dt}\omega + r^2\frac{d\omega}{dt}\right] \nonumber \\
&=m\left[2r\frac{dr}{dt}\omega + r^2\frac{d\omega}{dt}\right] \nonumber \\
&=mr\left[2\frac{dr}{dt}\omega + r\frac{d\omega}{dt}\right]
\end{align}
The part inside the brackets is the angular acceleration which for a central force is zero\footnote{see equation \ref{eq:angularAcceleration}} and so
\begin{align}
\frac{dl}{dt}&=mr \times 0 \nonumber \\
&=0
\end{align}
Which means that the angular momentum is constant.